% p163_1; source PDF page 163; manually verified editable table.
\begin{tikzpicture}[x=1cm,y=-1cm,inner sep=0pt,outer sep=0pt]
\path[use as bounding box] (0,0) rectangle (16.000,8.315);
\draw[line width=.4pt] (0.000,0.000) rectangle (2.200,1.205);
\node[align=center,text width=1.940cm,font={\linespread{1}\fontsize{10}{13.50}\selectfont\bfseries}] at (1.100,0.603) {作为义务人（甲）};
\draw[line width=.4pt] (2.200,0.000) rectangle (4.600,1.205);
\node[align=center,text width=2.140cm,font={\linespread{1}\fontsize{10}{13.50}\selectfont\bfseries}] at (3.400,0.603) {非作为义务人（乙）};
\draw[line width=.4pt] (4.600,0.000) rectangle (10.200,1.205);
\node[align=center,text width=5.340cm,font={\linespread{1}\fontsize{10}{13.50}\selectfont\bfseries}] at (7.400,0.603) {范例};
\draw[line width=.4pt] (10.200,0.000) rectangle (16.000,1.205);
\node[align=center,text width=5.540cm,font={\linespread{1}\fontsize{10}{13.50}\selectfont\bfseries}] at (13.100,0.603) {结论};
\draw[line width=.4pt] (0.000,1.205) rectangle (2.200,3.510);
\node[align=center,text width=1.940cm,font={\linespread{1}\fontsize{10}{13.50}\selectfont}] at (1.100,2.357) {帮助};
\draw[line width=.4pt] (2.200,1.205) rectangle (4.600,3.510);
\node[align=center,text width=2.140cm,font={\linespread{1}\fontsize{10}{13.50}\selectfont}] at (3.400,2.357) {实行};
\draw[line width=.4pt] (4.600,1.205) rectangle (10.200,3.510);
\node[align=left,text width=5.340cm,font={\linespread{1}\fontsize{10}{13.50}\selectfont}] at (7.400,2.357) {乙用欺诈合同欲诈骗丙公司，让丙公司的法务甲睁只眼、闭只眼。甲答应照办};
\draw[line width=.4pt] (10.200,1.205) rectangle (16.000,3.510);
\node[align=left,text width=5.540cm,font={\linespread{1}\fontsize{10}{13.50}\selectfont}] at (13.100,2.357) {乙构成诈骗罪的实行犯，甲构成不作为的帮助犯};
\draw[line width=.4pt] (0.000,3.510) rectangle (4.600,8.315);
\node[align=left,text width=4.340cm,font={\linespread{1}\fontsize{10}{13.50}\selectfont}] at (2.300,5.912) {作为义务的来源根据就是总论的“实质二分说”。常考的是先行行为。判断标准：先行行为若导致他人陷入危险状态，则负有保护义务（保护人身法益和财产法益的义务）};
\draw[line width=.4pt] (4.600,3.510) rectangle (10.200,8.315);
\node[align=left,text width=5.340cm,font={\linespread{1}\fontsize{10}{13.50}\selectfont}] at (7.400,5.912) {甲乙欲实施故意伤害罪，共同将丙女打晕。乙又要强奸丙或拿财物。甲看到，不阻止};
\draw[line width=.4pt] (10.200,3.510) rectangle (16.000,8.315);
\node[align=left,text width=5.540cm,font={\linespread{1}\fontsize{10}{13.50}\selectfont}] at (13.100,5.912) {甲构成强奸罪的不作为的帮助犯、盗窃罪的不作为的帮助犯};
\end{tikzpicture}
